\chapter{Abstract}
According to \ac{PCB}s and their digital circuits which are getting more complex all the time there is the requirement of being able to determine the test coverage of it. In dependency on these determined values then automatic boundary scan test chains could be performed or the need of changing something on the layout will be visible. This thesis describes the way of the conceptual design and the realization of a graphical tool for loading, showing different digital circuit diagrams projects which need to be in the \ac{ODB}++ format and the determination of test coverage parameters as well as the creation of \ac{SVF} files for being able to run the boundary scan test chains on these special devices which contains at least one \ac{JTAG} device. Furthermore there will be described detailed information about the frameworks, programming patterns and best practices which were used for creating this \ac{MVVM} C\# project and the project structure itself which was created for making sure to have the following key features available:
\begin{itemize}
	\item Opening any \ac{ODB}++ project by connecting to a stand alone \ac{GRPC} server and using the \ac{PCB}-Investigator API
	\item Exporting any opened projects to \ac{JSON} file for later usage and importing of that created files
	\item Creation and loading of project files which contains a manifest file for information about the project structure itself
	\item Modifying, exporting and importing of different kind of settings according the connection to the \ac{GRPC} server, the appearance and the test coverage values
	\item Determination of different test coverage values according to the boundary scan mechanism
	\item Creation of \ac{SVF} files for run the boundary scan test chain with a programmer from ProMik GmbH on a \ac{JTAG} containing device
\end{itemize}